\section{Introduction}\label{sec:1}

\subsection{The local-to-global architecture problem}\label{sec:11}

Each constituent of a software architecture is designed to be consistent
within its own local context. A service keeps its own data conventions, and
each individual handoff between modules satisfies its own contract.
Nevertheless, the system as a whole can exhibit semantic inconsistency. The
heart of this phenomenon is that the sum of local correctness does not
constitute global correctness.

This structure has the same shape as a classical problem type in
mathematics. Given a family of local data, where each local piece is
consistent and adjacent pieces are mutually consistent, does a global datum exist
that glues the whole together? Sheaf theory formalizes this question as the
sheaf condition, and cohomology measures the obstruction to gluing as a
class in $H^1$. This paper realizes this problem type on software
architecture and proves, as a comparison theorem, that the obstruction to
semantic repair and the obstruction of the geometry generated from the
architecture's simultaneous equation system coincide as the same
cohomology class.

\subsection{The SAGA research problem}\label{sec:12}

The research problem of SAGA (S\'emantique Architecturale, G\'eom\'etrie
Alg\'ebrique) is to set up the construction in the language of
\emph{repair} and the construction in the language of \emph{equations and
geometry} independently, and then to prove a comparison theorem connecting
their $H^1$ groups and residual classes (the two constructions are given in
Section~\ref{sec:4}). Once the comparison holds, semantic diagnoses
translate into geometric computations, and geometric computations translate
into semantic readings, in both directions.

Throughout this paper, the cohomology under discussion is the additive
\v{C}ech $H^1(\cU,-)$ relative to a selected monomorphic AAT cover $\cU$.
Identification with cover-independent sheaf cohomology is not among the
claims of this paper (\S\ref{sec:57}).

A remark on the name. The SAGA of this paper is unrelated to the Saga
pattern of compensating sequences known from distributed transactions
\citep{garciamolina1987sagas}. The name is an homage to Serre's GAGA
comparison theorem \citep{serre1956gaga}, which established the
correspondence between two geometries (\S\ref{sec:92}).

\subsection{The three layers of the result and the organization of the paper}\label{sec:13}

This paper presents the result in three layers. The mathematical layer is
the principal contribution of the paper: we construct the comparison
between semantic repair cohomology and equation-generated AAT \v{C}ech
cohomology, and prove the $H^1$ isomorphism, the residual class
correspondence, and the equivalence with global repair under the true
sheaf condition (Section~\ref{sec:5}). The Lean layer is the status at
release time: it records the formalized definitions, theorems, witnesses,
and proof chain at declaration granularity, fixing the correspondence with
the theorems of this paper and the axiom status (Section~\ref{sec:6}). The
ArchSig layer is a finite case study: ArchSig is a Rust command-line tool
that computes diagnoses from observations of the code and a selected
equation system; on a fixed microservice architecture, it shows in a
reproducible form how the mathematical objects are computed from finite
architecture evidence and how they change before and after repair
(Section~\ref{sec:7}).

In this paper, the \textbf{release identity} is the triple consisting of
the release tag of the Lean sources, the version of ArchSig, and the
digests of the input artifacts. The Lean sources and the ArchSig sources
are available in the public repository
\begin{center}
\href{https://github.com/iroha1203/AlgebraicArchitectureTheoryV2}{\nolinkurl{github.com/iroha1203/AlgebraicArchitectureTheoryV2}}.
\end{center}
The formalization status of Section~\ref{sec:6} and the reproduction
procedure of Section~\ref{sec:7} are stated against this fixed identifier.

The rest of the paper is organized as follows. Section~\ref{sec:2} presents
the AAT approach; Section~\ref{sec:3} the mathematical foundations needed
for SAGA; Section~\ref{sec:4} the independent construction of the two
complexes; Section~\ref{sec:5} the comparison theorem; Section~\ref{sec:6}
the Lean status; Section~\ref{sec:7} the ArchSig diagnosis of real code;
Section~\ref{sec:8} related work; Section~\ref{sec:9} the research
outlook; and Section~\ref{sec:10} the conclusion.
